% Generated draft: requires mathematical and source review.
\documentclass[11pt]{article}
\usepackage[margin=25mm]{geometry}
\usepackage{fontspec}
\setmainfont{Latin Modern Roman}
\usepackage{amsmath,amssymb,mathtools,unicode-math}
\setmathfont{Latin Modern Math}
\usepackage{enumitem,needspace,xurl,hyperref}
\setlength{\parindent}{0pt}
\setlength{\parskip}{5pt}
\setlength{\emergencystretch}{3em}
\newcommand{\N}{\mathbb{N}}
\newcommand{\Z}{\mathbb{Z}}
\newcommand{\Q}{\mathbb{Q}}
\newcommand{\R}{\mathbb{R}}
\newcommand{\C}{\mathbb{C}}
\newcommand{\F}{\mathbb{F}}
\newcommand{\A}{\mathbb{A}}
\newcommand{\PP}{\mathbb{P}}
\newcommand{\E}{\mathbb{E}}
\newcommand{\eps}{\varepsilon}
\newcommand{\dd}{\,\mathrm{d}}
\begin{document}
\section*{UnsolvedMath ID: 30002796}
\textbf{Source problem code:} catalogue entry
\subsection{Definitions and mathematical statement}
For terminology, let \(k\) be a field and let \(S=k[x_1,\ldots,x_N]\) have its standard grading. Given positive integers \(\lambda_1,\ldots,\lambda_r\) with \(\sum_{i=1}^r\lambda_i=d\), a degree-\(d\) form with the indicated factorization is a product \(F=F_1\cdots F_r\), where each \(F_i\in S_{\lambda_i}\) is nonzero. These are explanatory conventions: the supplied record does not specify the ground field, the relationship between \(N\) and \(n\), the allowed partitions \(\lambda\), or the precise definition of \(X_{n,\lambda}\).

For a standard graded Artinian \(k\)-algebra \(B=\bigoplus_{j\geq 0}B_j\), a weak Lefschetz element is a linear form \(\ell\in B_1\) such that
\[
\operatorname{rank}_k(\times\ell:B_j\longrightarrow B_{j+1})
=\min\{\dim_k B_j,\dim_k B_{j+1}\}
\qquad\text{for every }j\geq 0.
\]
A strong Lefschetz element satisfies
\[
\operatorname{rank}_k(\times\ell^q:B_j\longrightarrow B_{j+q})
=\min\{\dim_k B_j,\dim_k B_{j+q}\}
\qquad\text{for every }j\geq 0\text{ and }q\geq 1.
\]
These definitions distinguish possible meanings of the source's Lefschetz terminology; the record does not select either property or specify a successive Lefschetz condition on quotients.

\textbf{Source question, with unresolved qualifications.} For the Cohen--Macaulay algebras of codimension two described as arising from tangent-space ideals to \(X_{n,\lambda}\), do the Lefschetz elements required to determine improper intersections always exist? The record supplies neither the defining ideals and algebra construction nor the multiplication maps, degrees, number of elements, or successive quotients required by this question. If an Artinian reduction is intended, the reduction is also unspecified. Consequently, a fully quantified existence statement cannot be recovered from the supplied record.

The original formulation instead refers to algebras associated with ideals generated by powers of general linear forms. The record does not explain the relationship between this family and the tangent-space construction in the corrected formulation. Their equivalence is therefore not assumed.
\subsection{Short English statement}
The supplied question asks whether certain Cohen--Macaulay algebras associated with tangent spaces to varieties of reducible forms always admit enough Lefschetz elements to compute improper intersections. The algebra construction and the required Lefschetz conditions are missing, and the relationship to the original question about powers of general linear forms is unexplained. A precise conjecture requires these specifications.
\subsection{Sources}
\begin{itemize}
\item Local registry: \texttt{data/unsolvedmath/problems.json}; supplied complete source record, permanent UnsolvedMath ID \(30002796\), source problem code \texttt{catalogue entry}, title ``Lefschetz Elements for Powers of General Linear Forms.'' Collection: Oberwolfach Reports Open Problems.
\item Supplied report citation: \emph{Mini-Workshop: Ideals of Linear Subspaces, Their Symbolic Powers and Waring Problems} (2015), Oberwolfach Reports, pp.~489--532, DOI \texttt{10.4171/owr/2015/9}. \url{https://doi.org/10.4171/owr/2015/9}. The record attributes the problem to the workshop report; no individual proposer is supplied, and no problem-specific page or passage is identified.
\item Supplied literature source labeled ``Primary secant-variety result'': \url{https://doi.org/10.1016/j.jalgebra.2019.03.014}. The record supplies no authors, title, theorem number, or precise statement connecting this reference to the requested Lefschetz property.
\item Literature assessment dated exactly \texttt{2026-08-21}: the supplied curation labels the problem open and states that the relevant secant-variety formula and the assertion that the associated algebras always have enough Lefschetz elements remain conjectural beyond established special cases. This is supplied dated curation, not an independently verified assessment. The background explicitly calls for independent verification of time-sensitive status and cited claims.
\item The registry labels the statement \texttt{corrected\_verified} and reports that its mathematical content was preserved while notation, typography, or omitted definitions were repaired for publication. This supplied editorial designation does not resolve the missing formulation details or constitute independent mathematical verification here.
\end{itemize}
\end{document}
